% ==========================================================
% == Appendix                                              ==
% ==========================================================

% ----------------------------------------------------------
% Appendix A: Full training algorithm and hyperparameters
% ----------------------------------------------------------
\section{Full \method Training Algorithm}
\label{apdx:algorithm}

% [Appendix A 引子] 大意:给出完整伪代码 + 超参对照表,与 §4.5 的简化版对应。
% 中文对应内容:附录 A 给出 DACE 训练主循环的完整伪代码（Algorithm \ref{alg:dace-full}）,包括所有变量初始化、数据结构说明、边界情形与剪枝/回放阈值,并用超参数表 Table \ref{tab:hyperparams} 与 §5.1 及实现代码仓库一一对应。
Algorithm~\ref{alg:dace-full} gives the full pseudocode of \method's training loop, consistent with Algorithm~\ref{alg:dace-short} in the main paper but exposing all data structures, boundary cases, and hyperparameter defaults. Table~\ref{tab:hyperparams} lists the corresponding hyperparameter values and directly maps each to its location in the open-source implementation.

\begin{algorithm}[h]
  \caption{\method training loop (full)}
  \label{alg:dace-full}
  \begin{algorithmic}[1]
    \State \textbf{Input:} attacker $\piA$ (SFT-warm-started), defender $\piD$ (base instruct), seed set $\mathcal{Q}$
    \State \textbf{Hyperparams:} rounds $K$, per-stage steps $T$, group size $G$, batch sizes $B_{\mathrm{train}}, B_{\mathrm{replay}}$, decay $\gamma$, Beta prior $(\alpha,\beta)$, prune $(\epsilon, n_{\min})$, format coef, $\lambda_{\mathrm{succ}}=1, \lambda_{\mathrm{fail}}=0.5$, smoothing $\varepsilon$
    \State \textbf{State:} archive $\archive\leftarrow\emptyset$; strategy-count histogram $\mathbf{n}\in\mathbb{N}^{|\strategyspace|}\leftarrow\mathbf{0}$
    \For{round $k=1,\dots,K$}
      \For{each $x_i\in\archive$} \Comment{time decay (Eq.~\ref{eq:time-decay})}
        \State $s_i \leftarrow \gamma\,s_i;\quad f_i \leftarrow \gamma\,f_i$
      \EndFor
      \Statex \hspace*{-1.1em}\textit{\small \# ---- Stage 1: attacker RL (frozen $\piD$) ----}
      \For{step $t=1,\dots,T$}
        \State Sample seed batch $\{q_j\}_{j=1}^{B_{\mathrm{train}}}\sim\mathcal{Q}$; generate $G$ rollouts per $q_j$ from $\piA$
        \State For each rollout, extract $b(x)$ by regex + fuzzy substring; obtain $\piD(x)$ and judge verdicts
        \State Compute $R_{\mathrm{div}}(x)$ using Eq.~\ref{eq:coverage-gain} with shared denominator over the batch
        \State Compute $\lambda(x)$ (Eq.~\ref{eq:lambda}); form $R_A$ (Eq.~\ref{eq:reward-A})
        \State Update $\piA$ by GRPO on $R_A$ within each group
        \State For every successful $x$ not in $\archive$: insert $(x, b(x), s\!=\!1, f\!=\!0)$; increment $\mathbf{n}[b(x)]$
      \EndFor
      \Statex \hspace*{-1.1em}\textit{\small \# ---- Stage 2: defender RL (frozen $\piA$) ----}
      \For{step $t=1,\dots,T$}
        \State Generate $B_{\mathrm{train}}$ fresh attacks from $\piA$ on a seed batch
        \State For each $x_i\in\archive$: draw $\tilde p_i\sim\mathrm{Beta}(\alpha+s_i,\beta+f_i)$; let $\mathcal{D}_{\mathrm{replay}}\leftarrow\operatorname{Top-}B_{\mathrm{replay}}(\{\tilde p_i\})$
        \State Defender rolls out $G$ times per attack in $\mathrm{fresh}\cup\mathcal{D}_{\mathrm{replay}}$; compute $R_D$
        \State Update $\piD$ by GRPO on $R_D$ with the mixed mini-batch
        \For{each $x_i\in\mathcal{D}_{\mathrm{replay}}$}
          \State Apply Bayesian update Eq.~\ref{eq:beta-update} from the defender's verdict
          \If{$\frac{\alpha+s_i}{\alpha+\beta+s_i+f_i}<\epsilon$ \textbf{and} $s_i+f_i>n_{\min}$}
            \State Remove $x_i$ from $\archive$; decrement $\mathbf{n}[b(x_i)]$
          \EndIf
        \EndFor
        \State For every new success $x$ in fresh batch not already in $\archive$: insert entry; increment $\mathbf{n}$
      \EndFor
    \EndFor
    \State \Return $\piA,\piD,\archive$
  \end{algorithmic}
\end{algorithm}

% [Table A hyperparams] 占位超参表
\begin{table}[h]
  \centering
  \small
  \caption{\textbf{\method training hyperparameters.} Values shown are defaults used in all main-text experiments unless stated otherwise. ``Code path'' lists the file/function inside the released implementation where each hyperparameter is consumed.}
  \label{tab:hyperparams}
  \setlength{\tabcolsep}{4pt}
  \begin{tabular}{llll}
    \toprule
    \textbf{Group} & \textbf{Symbol} & \textbf{Default} & \textbf{Code path} \\
    \midrule
    Training       & $K$ rounds, $T$ steps, group $G$ & --, --, -- & \texttt{trainer.ppo.ray\_trainer} \\
    Batch          & $B_{\mathrm{train}}$, $B_{\mathrm{replay}}$ & --, -- & \texttt{trainer.ppo.ray\_trainer} \\
    Coverage       & smoothing $\varepsilon$, clip\_negative & $10^{-8}$, true & \texttt{archive\_pool.compute\_div} \\
    Diversity      & $\lambda_{\mathrm{succ}}, \lambda_{\mathrm{fail}}$ & $1.0, 0.5$ & \texttt{archive\_pool.compute\_div} \\
    Replay         & $\gamma$ decay & $0.97$ & \texttt{archive\_pool.time\_decay} \\
    Posterior      & $\alpha, \beta$ & $1.0, 1.0$ & \texttt{archive\_pool.\_post\_mean} \\
    Prune          & $\epsilon, n_{\min}$ & $0.05, 8$ & \texttt{archive\_pool.prune} \\
    Archive        & max\_pool\_size & $4000$ & \texttt{archive\_pool.add\_entry} \\
    \bottomrule
  \end{tabular}
\end{table}

% ----------------------------------------------------------
% Appendix B: Strategy space design
% ----------------------------------------------------------
\section{Strategy Space Design}
\label{apdx:space-design}

% [段 B.1 · 风险类别来源和剪裁证据] 大意:给出 12 类风险的设计理由,附统计证据。
% 中文对应内容:12 个风险类别来自 Llama-Guard-4 官方 14 类 taxonomy。我们对两类进行了剪裁：Sexual Content（原 S12）在 43.8k 条 SFT 蒸馏样本中频率仅 0.80%，且其 benign prompt 被 Guard 误判为 unsafe 的比例高达 96.3%（远高于其他类别），作为 diversity 信号不可靠；Code Interpreter Abuse（原 S14）在 SFT 数据中频率仅 0.11%，概念上指向工具调用与代码执行，text-only defender 下没有对应表面，无法构成真实威胁。剪裁后剩余 12 类对应 Llama-Guard-4 的 S1-S11 + S13。Table \ref{tab:risk-map} 给出映射;Figure \ref{fig:pruning-evidence} 展示剪裁前 14 类在 SFT + archive pool 中的频率与 benign mis-classification rate 的数据证据。
The 12 risk categories $\riskspace$ are derived from the Llama-Guard-4 taxonomy~\citep{llamaguard4_2026}. We prune two of its original 14 categories: \textit{Sexual Content}, whose benign prompts are mis-classified as unsafe by Llama-Guard-4 at a rate of $96.3\%$ on our distilled data (making the corresponding column of the diversity signal unreliable), and \textit{Code Interpreter Abuse}, whose concept presupposes tool-execution surfaces absent from text-only defenders, yielding only $0.11\%$ of SFT distillation samples. The remaining 12 risk categories correspond to Llama-Guard-4 codes S1--S11 and S13. Table~\ref{tab:risk-map} lists the mapping; Figure~\ref{fig:pruning-evidence} reports the supporting frequency and mis-classification statistics.

\begin{table}[h]
  \centering
  \small
  \caption{Mapping of \method's 12 risk categories to Llama-Guard-4 codes.}
  \label{tab:risk-map}
  \setlength{\tabcolsep}{4pt}
  \begin{tabular}{l l l}
    \toprule
    \textbf{\method Index} & \textbf{LG4 Code} & \textbf{Name} \\
    \midrule
    S1  & S1  & Violent Crimes \\
    S2  & S2  & Non-Violent Crimes \\
    S3  & S3  & Sex-Related Crimes \\
    S4  & S4  & Child Sexual Exploitation \\
    S5  & S5  & Defamation \\
    S6  & S6  & Specialized Advice \\
    S7  & S7  & Privacy \\
    S8  & S8  & Intellectual Property \\
    S9  & S9  & Indiscriminate Weapons \\
    S10 & S10 & Hate \\
    S11 & S11 & Suicide \& Self-Harm \\
    S12 & S13 & Elections \\
    --  & S12 & \textit{Sexual Content (pruned)} \\
    --  & S14 & \textit{Code Interpreter Abuse (pruned)} \\
    \bottomrule
  \end{tabular}
\end{table}

\begin{figure}[h]
  \centering
  \fbox{\rule[-2.0cm]{0pt}{4.0cm}\rule[-2.0cm]{0.95\linewidth}{0pt}}
  \caption{\textbf{Empirical evidence for the 12$\times$10 cut.} Per-category frequency in the SFT distillation set (left), archive pool at step 300 (middle), and benign mis-classification rate on Llama-Guard-4 (right). S12 Sexual Content and S14 Code Interpreter Abuse are outliers on all three axes; the pruned grid concentrates supervision on the 12 reliably covered cells.}
  \label{fig:pruning-evidence}
\end{figure}

% [段 B.2 · 攻击风格来源] 大意:10 攻击风格取自 Rainbow Teaming; 逐条定义;说明为何不再扩展（训练数据覆盖完整性 + 与 LG4 风险正交）。
% 中文对应内容:10 种攻击风格直接采用 Rainbow Teaming 的原始定义,它们覆盖了实践中最常见的语言层表面变化,且与 LG4 的语义风险分类大致正交。Table \ref{tab:style-def} 给出每一种风格的定义与典型例子。
\paragraph{Attack styles.} The 10 attack styles in $\stylespace$ are adopted without modification from Rainbow Teaming~\citep{samvelyan2024rainbow}: \textit{Slang, Technical Terms, Role Play, Authority Manipulation, Misspellings, Word Play, Emotional Manipulation, Hypotheticals, Historical Scenario, Uncommon Dialects}. Table~\ref{tab:style-def} lists per-style definitions and canonical examples.

\begin{table}[h]
  \centering
  \small
  \caption{Definitions of the 10 attack styles in $\stylespace$, following Rainbow Teaming.}
  \label{tab:style-def}
  \setlength{\tabcolsep}{4pt}
  \begin{tabular}{l l}
    \toprule
    \textbf{Style} & \textbf{Short Description} \\
    \midrule
    Slang                  & Colloquial or internet vernacular phrasing \\
    Technical Terms        & Jargon/specialist framing that obscures intent \\
    Role Play              & Fictional persona / character-driven context \\
    Authority Manipulation & Impersonating or invoking an authoritative role \\
    Misspellings           & Deliberate typos / orthographic perturbations \\
    Word Play              & Puns, homophones, or obfuscating word choice \\
    Emotional Manipulation & Appeals to urgency, sympathy, or distress \\
    Hypotheticals          & Counterfactual / ``what-if'' framings \\
    Historical Scenario    & Set in a past era or fictional historical context \\
    Uncommon Dialects      & Non-standard dialectal or multilingual register \\
    \bottomrule
  \end{tabular}
\end{table}

% [段 B.3 · SFT 热启动 + benign 模板] 大意:说明 Gemini-2.5-Pro 蒸馏 + 3 段式 CoT + v4 benign 模板修复,保证 attacker 输出格式遵从与低 benign 污染率。
% 中文对应内容:attacker 使用 Gemini-2.5-Pro 蒸馏约 43K 条三段式 CoT 样本进行 SFT 热启动,其中 benign 与 harmful 源 1:1 均衡。我们在 v4 版本中显式修复了 benign 模板:早期版本（v3)的 benign 模板混合了"欺骗 defender 生成 harmful 响应"与"保持 benign 语义"两个相互矛盾的目标,导致 43% 的 benign 样本在 Guard 上被判定为 unsafe。v4 把 benign 模板改为"在保持 benign 语义的前提下诱导 defender 过拒"的清晰目标,使 benign Guard not-safe 率降至 ≤ 10%。这是我们 benign-side 评估的数据洁净度基础。
The attacker is warm-started from Qwen2.5-7B-Instruct by SFT on $\sim$43K three-part CoT rewrites distilled from Gemini-2.5-Pro~\citep{comanici2025gemini}, balanced 1:1 between benign and harmful seeds. The benign-seed prompt template underwent an important revision: earlier versions mixed the goals of ``trick the defender into producing a harmful response'' and ``keep the rewrite benign,'' which were mutually inconsistent and led to $43\%$ of benign samples being flagged as unsafe by Llama-Guard-4. The revised template re-frames the benign-seed task as ``induce defender over-refusal while preserving benign intent,'' bringing the Guard ``not-safe'' rate on benign rewrites to $\leq 10\%$. We report full prompt templates (attacker, defender, judge) in Appendix~\ref{apdx:prompts}.

% ----------------------------------------------------------
% Appendix C: Prompts and Templates
% ----------------------------------------------------------
\section{Prompt Templates}
\label{apdx:prompts}

% [段 C 引子] 大意:给出三套 prompt 模板的位置与说明（详细内容可在 camera-ready 中完整列出）。
% 中文对应内容:本附录给出 attacker、defender 与 judge 三套 prompt 模板的完整内容。受篇幅限制,每套模板仅列关键段,完整版本随代码与数据释出。
This appendix lists the three prompt templates used throughout \method:
\paragraph{Attacker template.} A system message that (i) describes the full $12\times10$ strategy space with per-cell definitions, (ii) requires a three-part \texttt{<think>/<strategy>/<answer>} output format, and (iii) specifies the canonical strategy declaration syntax.
\paragraph{Defender template.} A standard safety-aware assistant system prompt used identically across all co-evolution baselines to avoid confounds.
\paragraph{Judge template.} A structured assessment prompt that elicits (a) a harmfulness label over LG4 S1--S11, S13, (b) a refusal label, and (c) a strategy-descriptor classification (used only for sanity checks of regex extraction). Full prompt text is released with the code repository.

% ----------------------------------------------------------
% Appendix D: Training setup
% ----------------------------------------------------------
\section{Training Setup}
\label{apdx:training-setup}

% [段 D] 大意:记录训练硬件、RL 框架、步数、curricula、随机种子协议等。
% 中文对应内容:附录 D 汇总 DACE 与各 baseline 的训练设置,包括硬件（GPU 数、VRAM）、RL 框架版本（veRL / GRPO 超参）、SFT 阶段使用的蒸馏模型版本、随机种子协议（三种子平均报告）与 checkpoint 保存策略。
\method and all co-evolution baselines are trained under identical compute budgets on an 8$\times$H100 node using veRL~\citep{sheng2025verl}. Unless otherwise stated, every reported number is the mean across three independent seeds; per-seed variance is reported in Appendix~\ref{apdx:expanded-results}. Checkpoints are saved at the end of every round and used both for co-evolution heatmaps (Figure~\ref{fig:coevolution-matrix}) and forgetting curves (Figure~\ref{fig:forgetting}).

% ----------------------------------------------------------
% Appendix E: Benchmark protocols
% ----------------------------------------------------------
\section{Benchmark Protocols}
\label{apdx:bench-protocol}

% [段 E] 大意:逐 benchmark 说明我们如何评测,特别是 MAJIC / Auto-RT / X-Teaming 这类需要特殊接入的 OOD attacker。
% 中文对应内容:我们对每个 benchmark 固定统一协议：safety 测评使用 GPT-4o 作为独立 judge,over-refusal 使用各 benchmark 自带判官规则;MAJIC 使用其作者释出的 attack pipeline 与默认超参直接攻击 defender API；Auto-RT 以 adapter 形式接入 defender 做 RL 迭代；X-Teaming 采用多轮 agent 化攻击并统计任一轮突破即算成功。
For each benchmark we fix a single evaluation protocol used identically across all methods. Static safety suites (WG-Test, WJB adv/benign, HarmBench adv/van, DAN, XSTest, OR-Bench, StrongREJECT) are scored using both Qwen3Guard (for training) and GPT-4o (for final reporting). OOD attacker evaluations use the attackers' own published pipelines against an API-wrapped defender: PAIR/TAP/AutoDAN/AutoDAN-Turbo use their default query budgets; MAJIC runs under its default Markovian composition schedule; Auto-RT is attached as an adversarial RL adapter; X-Teaming uses multi-turn agentic rollouts where success at any turn counts as a break. General-capability benchmarks follow official protocols (IFEval, ARC-C, GPQA via \textsc{lm-eval-harness}; MMLU via 5-shot; AlpacaEval~2 with length-controlled winrate).

% ----------------------------------------------------------
% Appendix F: Expanded results
% ----------------------------------------------------------
\section{Expanded Results}
\label{apdx:expanded-results}

% [段 F] 大意:多 backbone、多 judge 交叉验证结果,与 main text 表 1-3 对齐,附录加上 Llama3.1-8B-Instruct 与 Qwen2.5-14B-Instruct 两个扩展 backbone。
% 中文对应内容:附录 F 给出 §5 主结果在 Llama3.1-8B-Instruct 与 Qwen2.5-14B-Instruct 两个扩展 backbone 下的完整版本;每个 backbone 下重复 RQ1-RQ3 的全部 benchmark;此外给出 Qwen3Guard vs GPT-4o vs Llama-Guard-4 三 judge 交叉验证的一致性分析,以排除 judge 引入的系统性偏置。
This appendix repeats the main-text Tables~\ref{tab:rq1-safety}, \ref{tab:rq2-capability}, and \ref{tab:rq3-ood} on two additional defender backbones (Llama3.1-8B-Instruct and Qwen2.5-14B-Instruct), and contains a three-judge cross-validation (Qwen3Guard, GPT-4o, Llama-Guard-4) that confirms the ranking of methods is robust to judge choice. Detailed tables are placeholders and will be populated once evaluation completes.

% ----------------------------------------------------------
% Appendix G: Training dynamics
% ----------------------------------------------------------
\section{Training Dynamics}
\label{apdx:training-dynamics}

% [段 G] 大意:给出 DACE 训练过程中 pool size、entropy、coverage reward、forgetting 等训练指标的演化曲线,与 main text 的 forgetting 图 5 互补。
% 中文对应内容:附录 G 展示 DACE 训练过程中的关键动力学曲线: pool size 从 0 平滑增长到饱和阈值；策略分布熵 H(p_{π_A,t}) 从早期低值逐步上升接近 log 120;归一化覆盖增益奖励 mean ± std 在全训练周期保持 O(1) 量级无 vanishing；time decay 使 defender stage 的贝叶斯 posterior mean 在每轮开始时轻微回落;forgetting curve 的放大版进一步揭示各 replay 策略在训练轴不同位置的效果差异。
Figure~\ref{fig:training-dynamics} plots the full set of training-time dynamics of \method: archive size, strategy-distribution entropy $H(\mathbf{p}_{\piA,t})$, the mean and standard deviation of the normalized coverage reward, posterior means before/after time decay at each round boundary, and the per-time-slice forgetting curve (an enlarged version of Figure~\ref{fig:forgetting}). These confirm that (i) the coverage reward never vanishes, (ii) strategy-distribution entropy monotonically approaches $\log|\strategyspace|$, and (iii) Bayesian posterior decay properly tracks the defender's capability drift.

\begin{figure}[h]
  \centering
  \fbox{\rule[-2.5cm]{0pt}{5.0cm}\rule[-2.5cm]{0.95\linewidth}{0pt}}
  \caption{Training dynamics of \method (placeholder; curves rendered from training logs).}
  \label{fig:training-dynamics}
\end{figure}

% ----------------------------------------------------------
% Appendix H: Attack patterns qualitative analysis
% ----------------------------------------------------------
\section{Attack Patterns and Combinatorial Emergence}
\label{apdx:attack-patterns}

% [段 H] 大意:给出若干 sanitized 组合攻击案例（与 Figure 4 对应）,每例配有策略描述符和简要解读。
% 中文对应内容:附录 H 选取 DACE 训练后期涌现的若干组合攻击案例,按策略类别 × 风格的方式呈现(prompt 内容经伦理清洗,保留结构性特征)；每例标注所选 (s,c) 对与最终 ASR,展示 DACE attacker 如何在单个 rewrite 中 seamless 地融合多类风险和多种风格。案例集在附录 table 中全部给出,并在 caption 中注明"攻击细节已经过 sanitization,不直接暴露可操作内容"。
We present eight sanitized examples of combinatorial rewrites produced by \method in mid-to-late training (Table~\ref{tab:combinatorial}). Each example is annotated with the self-declared strategy descriptor $b(x)=(s,c)$, the strategies \emph{actually} engaged by the text (judged post-hoc), and the defender's final outcome. All operational details have been redacted; we retain only the structural composition pattern. The qualitative pattern---seamless composition of risk categories with stylistic framings---is absent from MAGIC and TriPlay-RL outputs under the same seed prompts.

\begin{table}[h]
  \centering
  \small
  \caption{Combinatorial attack examples from \method (placeholder; contents to be added post-sanitization). Annotations: declared $(s,c)$, secondary strategies engaged, judge outcome.}
  \label{tab:combinatorial}
  \begin{tabular}{p{0.08\linewidth} p{0.28\linewidth} p{0.28\linewidth} p{0.24\linewidth}}
    \toprule
    \textbf{\#} & \textbf{Declared $(s,c)$} & \textbf{Secondary strategies engaged} & \textbf{Outcome} \\
    \midrule
    1 & -- & -- & -- \\
    2 & -- & -- & -- \\
    \ldots & -- & -- & -- \\
    \bottomrule
  \end{tabular}
\end{table}

% ----------------------------------------------------------
% Appendix I: Full Related Work
% ----------------------------------------------------------
\section{Extended Related Work}
\label{apdx:related-work-full}

% [段 I.1] LLM 安全对齐（背景）
% 中文对应内容:早期 LLM 安全对齐以 RLHF、Constitutional AI 与 Safe RLHF 等方法为主,重点是用人类/AI 反馈对齐一次性分布;随 jailbreak 攻击出现,静态的 red-teaming 数据集（WildGuard、HarmBench、Sorry-Bench、OR-Bench、StrongREJECT、XSTest 等)被用来扩充对齐数据;但这些方法均假设攻击分布接近某个"固定靶"。
\paragraph{Static safety alignment.} Foundational alignment methods---RLHF~\citep{ouyang2022training}, Constitutional AI~\citep{bai2022constitutional}, Safe RLHF~\citep{dai2023safe}, DPO~\citep{rafailov2023dpo}---all optimize against a fixed distribution of preferences or demonstrations. Static red-teaming data then extends this distribution~\citep{jiang2024wildteaming, mazeika2024harmbench, rottger2023xstest, cui2024orbench, souly2024strongreject, xie2024sorry}. The moving-target problem~\citep{liu2025chasing} motivates the co-evolutionary direction we build on.

% [段 I.2] 自动红队与 jailbreak 全景
% 中文对应内容:自动红队包括梯度/优化类（GCG、AutoDAN),LLM-as-attacker 类（PAIR、TAP）、多轮/agentic 类（X-Teaming、MAJIC）,以及最新的 Markovian 组合式(MAJIC)。这些工作主要关心"如何让一个固定 defender 被攻破",与 DACE 的"让 attacker 的策略分布结构化覆盖"互补。
\paragraph{Automated red-teaming.} A rich literature develops automated attackers of varying paradigms: gradient/optimization (GCG~\citep{zou2023universal}, AutoDAN~\citep{liu2023autodan}), LLM-as-attacker iterative rewrites (PAIR~\citep{chao2023pair}, TAP~\citep{mehrotra2024tap}, AutoDAN-Turbo~\citep{liu2024autodan}), multi-turn or agentic attackers (X-Teaming~\citep{rahman2025xteaming}, MAJIC~\citep{qi2025majic}), and scenario-shift attacks (FlipAttack~\citep{liu2024flipattack}). These works typically treat the defender as fixed; \method instead integrates a \emph{diverse} attacker inside a co-training loop.

% [段 I.3] 协同演化详细对比
% 中文对应内容:协同演化先验工作在结构上可以进一步细分: 零和自博弈（Self-RedTeam)、非对称序贯博弈（MAGIC)、多智能体（AdvEvo-MARL），tri-role（TriPlay-RL)、树组双感知（ACE-Safety)、自演化（SEAS)、Stackelberg 博弈、非合作博弈（AdvGame)、多轮自动红队（MART）等。表 \ref{tab:coevo-compare} 给出结构化对比,说明 DACE 是首个在协同演化中同时处理策略空间覆盖与连续不确定性记忆的方法。
\paragraph{Co-evolutionary alignment: detailed comparison.} Table~\ref{tab:coevo-compare} summarizes structural differences between \method and prior co-evolutionary frameworks. \method is the first to (i) introduce a strategy-indexed coverage objective with an $\mathcal{O}(1)$-signal KL-contraction reward, and (ii) maintain threat memory via continuous Beta--Bernoulli posteriors with time decay and Thompson sampling.

\begin{table}[h]
  \centering
  \small
  \caption{Structural comparison of co-evolutionary safety-alignment frameworks.}
  \label{tab:coevo-compare}
  \setlength{\tabcolsep}{3pt}
  \begin{tabular}{l c c c c c c}
    \toprule
    \textbf{Method} & \textbf{Game form} & \textbf{Explicit strategy space} & \textbf{Coverage reward} & \textbf{Replay pool} & \textbf{Threat model} & \textbf{Time decay} \\
    \midrule
    Self-RedTeam~\citep{liu2025chasing} & zero-sum self-play & no  & no  & no  & -- & no \\
    MAGIC~\citep{magic2025}             & asym.\ sequential  & no  & no  & no  & -- & no \\
    AdvEvo-MARL~\citep{pan2025advevo}   & multi-agent        & no  & no  & seed pool & -- & no \\
    SSP~\citep{ssp2026}                 & self-play          & no  & no  & yes & discrete score & no \\
    TriPlay-RL~\citep{triplayrl2026}    & tri-role self-play & no  & Self-BLEU+cos (textual) & no & -- & no \\
    ACE-Safety~\citep{acesafety2025}    & tree-group MCTS    & no  & no  & no  & -- & no \\
    SEAS~\citep{seas2025}               & self-evolve        & no  & no  & no  & -- & no \\
    \rowcolor{cyan!8}\textbf{\method (ours)} & asym.\ sequential & \textbf{yes $12\!\times\!10$} & \textbf{normalized KL contraction} & \textbf{yes} & \textbf{Beta--Bernoulli continuous} & \textbf{yes} \\
    \bottomrule
  \end{tabular}
\end{table}

% [段 I.4] Diversity in Red-teaming 扩展 + continual learning analogy
% 中文对应内容:除了 §2.3 提到的 QD 流派与 RL-intrinsic 多样性流派,相关领域还包括持续学习中的经验回放（experience replay)与非平稳多臂赌博机（non-stationary bandit)文献;我们的 Beta-Bernoulli 回放在形式上最接近后者的 Thompson Sampling 方案,但针对的是防御者能力漂移下的攻击风险估计问题。
\paragraph{Experience replay and non-stationary bandits.} Our Bayesian replay mechanism draws on two complementary traditions: experience replay in continual RL~\citep{lin2020extrapolation} and non-stationary Thompson sampling in bandit theory~\citep{chapelle2011thompson}. The unique adaptation in \method is to treat each archived attack as a Bernoulli arm whose latent threat probability drifts with the defender, and to propagate drift through a time-decayed Beta posterior so that the replay selection step doubles as a risk re-estimation step.

% ----------------------------------------------------------
% Appendix J: Limitations, Broader Impact, Ethics
% ----------------------------------------------------------
\section{Limitations, Broader Impact, and Ethical Considerations}
\label{apdx:limitations}

% [段 J.1 Limitations expanded] 大意:对 §6 的局限做扩展分析。
% 中文对应内容:首先,12×10 的策略空间仍然是人工设计且离散,未来可探索自动发现的行为空间（例如基于 embedding 聚类或无监督结构发现）。其次,组合涌现依赖 SFT 种子数据的质量与判官模型的稳定性;若上游存在偏置（如 Guard 对某些风险类别系统性低估）,DACE 的协同演化可能放大这些偏置。第三,本文仅在文本单轮场景验证;DACE 向多模态、多轮 agent 与工具使用场景的推广仍需系统实验。
We expand on the limitations briefly noted in Section~\ref{sec:conclusion}. \emph{Discretization.} The $12\times10$ grid, while empirically well-chosen, is still a coarse human-designed quantization of the threat space; automatic behavior-space discovery is a natural extension. \emph{Judge bias propagation.} \method's coverage and replay signals both consume a safety judge; any systematic bias of the judge (e.g., under-estimation of some risk categories) may be amplified by co-evolution. Multi-judge cross-validation (Appendix~\ref{apdx:expanded-results}) mitigates but does not eliminate this risk. \emph{Modality and interaction scope.} Extending \method to multimodal defenders, multi-turn dialogue, and tool-using agents requires redesigning the strategy grid and possibly generalizing the Bayesian posterior to richer outcome spaces.

% [段 J.2 Broader Impact]
% 中文对应内容:DACE 旨在强化 LLM 安全对齐;训练过程涉及生成更多样、更强的攻击样本,这些样本若被滥用可能提高对现有模型的攻击成功率。我们采取三类保障措施:(i) 只发布 defender 权重与训练代码,不发布 attacker 权重;(ii) archive 中的攻击 prompt 在释出时做 ethical sanitization,模糊可操作细节;(iii) 社区研究者访问完整 archive 需签署非滥用协议。
\paragraph{Broader impact and safeguards.} \method is designed to strengthen safety alignment, but the training process inevitably generates higher-quality, more diverse adversarial prompts that could, in principle, be misused. We adopt three safeguards: (i) the public release includes the defender weights, training code, and evaluation infrastructure, but \emph{not} the attacker weights; (ii) the released archive of successful attacks is ethically sanitized to preserve structural composition patterns while redacting operational details; (iii) access to the un-sanitized archive is gated behind a non-misuse agreement. On the positive side, by making strategy-space coverage and adversarial memory explicit, \method makes the resulting defenders easier to audit: reviewers can directly inspect which risk categories and attack styles the defender was trained against and how the threat posteriors evolved.

% [段 J.3 Ethical considerations]
% 中文对应内容:本研究在生成/评估阶段全部使用公开 benchmark 与合成 prompts,不涉及真实用户数据或敏感个人信息;判官模型与 attacker/defender 的训练均在机构内部计算集群完成。我们不鼓励将 DACE 的 attacker 直接用于任何未经授权的系统越狱。
\paragraph{Ethical considerations.} All experiments use publicly released benchmarks and synthetically generated prompts; no real user data or personally identifying information is involved. Training and evaluation are performed on institutional compute. \method's attacker is not to be used against production systems without explicit authorization.
